% ==========================================
% BAB II STUDI LITERATUR
% ==========================================
% Penanda komentar (tidak tampil di PDF):
%   % SUMBER-BELUM-ADA : klaim butuh sitasi yang belum dimiliki/dibaca.
% Sitasi tanpa penanda sudah dicek terhadap PDF.
% ==========================================
\chapter{STUDI LITERATUR}
\label{chap:studi-literatur}
Bab ini membahas petir dan parameter yang memengaruhinya, komputasi kuantum dan QNN, serta kesenjangan penelitian.

\section{Petir dan \textit{Ground Flash Density} (GFD)}

\subsection{Pembentukan Petir}
Petir terbentuk dari pemisahan muatan di zona fase campuran awan konvektif, yaitu lapisan yang mengandung air sangat dingin, es, dan \textit{graupel} \cite{shanTBDMLI}.
Kristal es yang bertumbukan dengan \textit{graupel} menjadi bermuatan positif dan terangkat oleh arus naik, sedangkan \textit{graupel} bermuatan negatif dan jatuh \cite{cheng2024MLLP}.
Karena itu, parameterisasi petir umumnya bertumpu pada ukuran badai, intensitas konveksi, dan jumlah hidrometeor es \cite{cheng2024MLLP}.
Petir terbagi menjadi petir awan ke tanah (CG) dan petir dalam awan (IC) \cite{shanTBDMLI}, dan penelitian ini hanya membahas petir CG.

\subsection{GFD dan Pengukurannya}
Densitas kilat adalah jumlah kilat per satuan luas per satuan waktu \cite{peterson2021LISOTD}, sedangkan GFD adalah densitas kilat yang hanya menghitung petir CG.
Sensor satelit seperti LIS dan OTD mendeteksi petir CG dan IC secara bersamaan sehingga tidak dapat langsung menghasilkan GFD \cite{peterson2021LISOTD}.
Jaringan sensor darat menentukan lokasi setiap sambaran, tetapi efisiensi deteksinya menurun seiring jarak dari pusat jaringan \cite{song2023LNAI}.
Penelitian ini menggunakan data sistem deteksi petir PLN Puslitbang untuk Jawa Barat dan arsip MERLIN untuk Florida \cite{ksc2026merlin}.
% SUMBER-BELUM-ADA: dokumentasi sistem deteksi PLN.
Klasifikasi CG/IC pada MERLIN belum dipastikan sehingga hasil domain subtropis bersifat sementara.
Penurunan efisiensi deteksi terhadap jarak juga tidak dapat dipisahkan dari pola klimatologis menggunakan data MERLIN saja.

\subsection{Karakteristik Petir Tropis dan Subtropis}
Kelembapan tinggi, arus naik yang kuat, serta aerosol dari industri dan garam laut memudahkan pembentukan awan kumulonimbus di Indonesia \cite{denov2022}.
Standar proteksi petir internasional terutama dikembangkan dari karakteristik petir subtropis, padahal arus puncak, frekuensi, dan GFD petir tropis berbeda \cite{denov2025RNFPA}.
Perbedaan lingkungan juga mengubah peran variabel, misalnya CAPE kurang berkorelasi dengan densitas kilat di tenggara Amerika Serikat karena CAPE di wilayah tersebut hampir selalu cukup untuk memicu badai \cite{cheng2024MLLP}.

\subsection{Parameter yang Memengaruhi GFD}
Konveksi dalam membutuhkan ketidakstabilan, kelembapan, dan pengangkatan \cite{haywardTBD}, sedangkan petir juga membutuhkan es di zona fase campuran.
Kandidat variabel dikelompokkan berdasarkan peran tersebut dan dirangkum beserta rujukannya pada Tabel \ref{tbl:tabelKandidatVariabel}.
Sebagian besar kandidat berasal dari ERA5 \cite{hersbach2020era5}, yaitu reanalisis per jam dengan resolusi 0,25° \cite{shanTBDMLI}, dan lima kandidat berasal dari NASA POWER \cite{nasapower}.
Karena ERA5 adalah reanalisis, model pada penelitian ini bersifat diagnostik dan kinerjanya merupakan batas atas kinerja operasional.

\input tables/tabelKandidatVariabel.tex

Ketidakstabilan diwakili oleh CAPE, CIN, indeks K (KX), dan indeks \textit{total totals} (TOTALX).
CAPE termasuk tiga prediktor terpenting kejadian petir, tetapi tidak lagi membedakan kejadian petir setelah mencapai nilai tinggi \cite{shanTBDMLI}.
CIN digunakan sebagai prediktor pada prakiraan petir \cite{sobash2025MLP}.
% SUMBER-BELUM-ADA: sumber rumus dan studi petir untuk KX dan TOTALX.

Kelembapan diwakili oleh RH2M, D2M, TCWV, VIMDF, CBH, dan variabel turunan \texttt{dewpoint\_depression} (T2M dikurangi D2M).
Kelembapan relatif dan titik embun permukaan digunakan pada model prediksi petir \cite{shanTBDMLI, sobash2025MLP}, sedangkan air mampu endap kolom total digunakan untuk mengenali pola sinoptik pemicu badai petir \cite{haywardTBD}.
Selisih suhu dan titik embun mewakili ketinggian kondensasi, yang menentukan kedalaman awan hangat dan jumlah tetes air yang mencapai zona fase campuran \cite{shanTBDMLI}.
% SUMBER-BELUM-ADA: VIMDF; hubungan selisih titik embun dengan ketinggian kondensasi.

Kondensat awan diwakili oleh TCLW, TCIW, VIIWD, VILWD, dan variabel turunan \texttt{cloud\_water} (TCLW ditambah TCIW).
Jumlah hidrometeor es merupakan komponen utama parameterisasi petir \cite{cheng2024MLLP}, dan zona fase campuran juga membutuhkan air sangat dingin \cite{shanTBDMLI}.
Karena itu, \texttt{cloud\_water} menjumlahkan kondensat cair dan es tanpa membedakan fase.
% SUMBER-BELUM-ADA: VIIWD dan VILWD.

Hujan diwakili oleh PRECTOTCORR dan laju hujan konvektif (CRR).
Presipitasi berkorelasi tinggi dengan densitas kilat, tetapi menjadi penduga yang buruk di wilayah dengan hujan non-konvektif \cite{cheng2024MLLP}, sehingga CRR juga digunakan.
Kondisi permukaan diwakili oleh PS, T2M, dan WS2M, yang bersama RH2M membentuk model empat variabel permukaan dengan kemampuan prediksi petir yang cukup baik \cite{shanTBDMLI}.

Lintang (\texttt{lat}) dan bujur (\texttt{lon}) digunakan sebagai prediktor statis \cite{sobash2025MLP}.
Jam dan hari dalam setahun bersifat siklis sehingga dikodekan dengan pasangan sinus dan kosinus (\texttt{hour\_sin}, \texttt{hour\_cos}, \texttt{doy\_sin}, \texttt{doy\_cos}) \cite{sobash2025MLP}.
Kosinus sudut zenit matahari (\texttt{cos\_sza}) menggabungkan kedua siklus tersebut dan dimotivasi oleh hubungan laju kilat dengan penyinaran matahari \cite{peterson2021LISOTD}.
Tidak ditemukan studi petir atau konveksi yang menggunakan \texttt{cos\_sza} sebagai prediktor sehingga variabel ini diuji tanpa preseden.

CIN dan CBH diambil tetapi tidak dimodelkan karena pada model ECMWF keduanya sengaja memuat nilai kosong \cite{ecmwf2015cy41r1}, misalnya CIN pada kolom tanpa dasar awan \cite{ecmwf2026era5doc}.
Nilai kosong tersebut menandakan besaran yang tidak terdefinisi, bukan data hilang.
\textit{Aerosol optical depth} tidak diambil karena NASA POWER hanya menyediakannya secara bulanan \cite{nasapower}, meskipun aerosol memengaruhi petir \cite{song2023LNAI}.

\subsection{Pendekatan Prediksi Petir}
Parameterisasi petir pada model cuaca dan iklim sangat beragam dan dapat menghasilkan proyeksi yang saling bertentangan \cite{kalashnikov2024PCG}.
Pembelajaran mesin menjadi alternatif, misalnya \textit{random forest} \cite{shanTBDMLI}, LightGBM \cite{song2023LNAI}, jaringan saraf \cite{sobash2025MLP}, dan \textit{convolutional neural network} \cite{kalashnikov2024PCG}.
Dua tantangan muncul berulang, yaitu sampel dengan petir jauh lebih sedikit daripada sampel tanpa petir \cite{song2023LNAI} dan variabel meteorologis yang saling berkorelasi \cite{cheng2024MLLP}.

\section{Komputasi Kuantum dan \textit{Quantum Neural Network} (QNN)}

\subsection{Dasar Komputasi Kuantum}
Komputer kuantum memproses informasi menggunakan \textit{qubit} yang dapat berada dalam superposisi dan saling terkait (\textit{entangled}).
Gerbang kuantum mengolah \textit{qubit}, dan pengukuran di akhir sirkuit menghasilkan bit klasik \cite{oliveira2024QNN}.
Perangkat kuantum saat ini tergolong \textit{Noisy Intermediate-Scale Quantum} (NISQ) yang belum dapat memberikan keunggulan kuantum sehingga banyak penelitian menggunakan simulasi sirkuit kuantum \cite{oliveira2024QNN}.

\subsection{QNN dan Alasan Pemilihannya}
QNN terdiri atas \textit{feature map} yang memetakan data ke keadaan kuantum, \textit{ansatz} berparameter yang dilatih oleh pengoptimal klasik, dan pengukuran \cite{oliveira2024QNN}.
Pada prakiraan iradiasi matahari, QNN memberikan hasil yang kompetitif untuk horizon 5 hingga 120 menit \cite{oliveira2024QNN}.
QSVR telah diterapkan pada data petir Jawa Barat dan Florida \cite{fadhil2025}, sedangkan QNN belum.
Karena itu, penelitian ini menggunakan QNN.
Perbandingan terstruktur antaralgoritma kuantum disajikan pada Bab III.

\subsection{Penyiapan Data untuk QNN}
Pada \textit{angle encoding}, setiap fitur menjadi sudut rotasi satu \textit{qubit} sehingga jumlah \textit{qubit} sama dengan jumlah prediktor \cite{oliveira2024QNN}.
Karena itu, fitur harus dipilih dan diskalakan ke rentang sudut rotasi.
Penelitian ini memilih fitur dengan \textit{minimum redundancy maximum relevance} (mRMR) berbasis \textit{mutual information} dan mengonfirmasinya dengan ablasi \textit{leave-one-out}.
% SUMBER-BELUM-ADA: mRMR (Peng, Long, & Ding 2005) dan rentang skala angle encoding.
Hasil ablasi perlu ditafsirkan hati-hati karena variabel meteorologis saling berkorelasi \cite{cheng2024MLLP}.
Target jumlah kilat per jam didominasi nilai nol sehingga dilatih dalam skala $\log(1+y)$ tanpa mengisi nilai nol dari metrik lain.

\section{Penelitian Terdahulu dan \textit{Research Gap}}
Penelitian prediksi petir berbasis pembelajaran mesin yang ditinjau seluruhnya dilakukan di Amerika Serikat \cite{cheng2024MLLP, shanTBDMLI, song2023LNAI, kalashnikov2024PCG, sobash2025MLP}.
Di bidang kuantum, QNN telah diterapkan pada prakiraan iradiasi matahari \cite{oliveira2024QNN}, dan QSVR telah diterapkan pada data petir Jawa Barat dan Florida \cite{fadhil2025}.
Dari penelitian tersebut, terdapat tiga kesenjangan.

Dari perspektif geografis, pengetahuan dan standar tentang petir terutama dikembangkan dari karakteristik petir subtropis, padahal petir tropis berbeda \cite{denov2025RNFPA}.
Model yang hanya dikembangkan dari data Amerika Serikat belum tentu berlaku di wilayah tropis seperti Jawa Barat.

Dari perspektif metodologis, hubungan variabel meteorologis dengan petir berbeda antarwilayah.
Presipitasi merupakan penduga yang baik di tenggara Amerika Serikat, tetapi buruk di pantai barat \cite{cheng2024MLLP}, dan keberhasilan indeks konveksi bervariasi menurut wilayah dan musim \cite{haywardTBD}.
Karena itu, model dan variabelnya perlu diuji pada dua rezim iklim yang berbeda, termasuk \texttt{cos\_sza} yang belum pernah diuji sebagai prediktor petir.

Dari perspektif teknis, perangkat NISQ belum dapat memberikan keunggulan kuantum \cite{oliveira2024QNN}.
Karena itu, QNN perlu dibandingkan secara langsung dengan model klasik pada data yang sama, bukan diasumsikan unggul.

Penelitian ini mengisi ketiga kesenjangan tersebut dengan membandingkan QNN variasional dan model klasik untuk prediksi jumlah kilat CG per jam di Jawa Barat dan Florida.
